\bigskip




% lecture 5, page 1
% \section*{Lecture 5 - Page 1}


\begin{center}
\fbox{\textbf{ALGEBRA}}\\[2pt]
\fbox{(Lecture 5)}
\end{center}
\begin{flushright}
\textit{1.XI.1995}
% \footnote{the exact date is unclear in the original -- it could be 1.XI, or the year 1994.}
\end{flushright}

\bigskip
\noindent\underline{\textbf{Lagrange's Theorem}}

\bigskip
\noindent Let $(G,\cdot)$ be a group. $H\leq G$.

\noindent If $x,y\in G$ \quad $x\equiv_L y \pmod H \overset{\text{def}}{\Longleftrightarrow} x^{-1}y\in H$.

\[
R_H^L = \{(x,y)\in G\times G \mid x^{-1}y\in H\} \subseteq G\times G.
\]

\noindent $R_H^L$ is an equivalence relation on $G$, compatible on the left with $\cdot$.

\[
xR_H^L y \Rightarrow zxR_H^L zy
\]

\noindent For $a\in G$, \ $[a]_{R_H^L} = aH = \{ah \mid h\in H\}$ = the left coset

\bigskip
\noindent $x\equiv_R y \pmod H \overset{\text{def}}{\Longleftrightarrow} xy^{-1}\in H$.

\[
R_H^R = \{(x,y)\in G\times G \mid xy^{-1}\in H\} = \text{an equivalence relation compatible on the right with } \cdot.
\]

\[
xR_H^R y \Rightarrow xzR_H^R yz.
\]

\[
[a]_{R_H^R} = Ha = \{ha \mid h\in H\}.
\]

\bigskip
\noindent Hence $T=\{a_i\}_{i\in I}$ = a left transversal of $H$ in $G$

\noindent = a complete system of representatives

\begin{enumerate}
\item[1)] $i\neq j \Rightarrow a_i \not\equiv_L a_j \pmod H$
\item[2)] $(\forall)\, x\in G$ \ $\exists\, i\in I$ such that $x\equiv_L a_i \pmod H$
\end{enumerate}

\bigskip
\noindent\underline{Notation.} $T^{-1} = \{a_i^{-1}\}_{i\in I}$

\noindent $T$ a left transversal $\Rightarrow$ $T^{-1}$ is a right transversal.

\[
T \xrightarrow{\ } T^{-1}, \qquad a_i \mapsto a_i^{-1} \text{ is a bijection} \ \Rightarrow \ \text{card}\, T = \text{card}\, T^{-1}
\]

\noindent = the index of $H$ in $G \overset{\text{not}}{=} [G:H]$

\noindent ord $G$ = card $G$ = $|G|$.

\bigskip
\noindent\fbox{Lagrange's Theorem} \quad If $(G,\cdot)$ is a finite group and $H$ is a subgroup of $G$,

\noindent then:

\[
\boxed{|G| = |H|\cdot [G:H]}
\]

\noindent or, equivalently: \quad $\big($ord $G$ = ord $H\cdot [G:H]\big)$

\bigskip
\noindent\underline{Application:} ord $S_n = n!$

\[
S_n = \{\sigma : \{1,2,\ldots,n\}\to\{1,2,\ldots,n\} \mid \sigma \text{ is bijective}\}.
\]

% \noindent [Notation:]\footnote{the connecting fragment here is very compressed in the original.}

\[
\sigma = \begin{pmatrix} 1 & 2 & \cdots & n \\ \sigma(1) & \sigma(2) & \cdots & \sigma(n) \end{pmatrix}
\]
